% ECONOMICS_PROBLEM: {"id": "E37-594", "title": "Recession and tail-risk forecasting", "jel_code": "E37", "source_id": "OP-0239"}
% !TeX program = xelatex
\documentclass[11pt,a4paper]{article}
\usepackage[margin=23mm]{geometry}
\usepackage{fontspec}
\setmainfont{Latin Modern Roman}
\usepackage{amsmath,amssymb,mathtools}
\usepackage{xcolor,enumitem,booktabs,longtable,array,xurl,hyperref}
\definecolor{muted}{HTML}{53636C}
\hypersetup{colorlinks=true,urlcolor=blue,linkcolor=blue}
\setlength{\parindent}{0pt}
\setlength{\parskip}{5pt}
\newcommand{\R}{\mathbb R}
\newcommand{\E}{\mathbb E}
\newcommand{\N}{\mathbb N}
\newcommand{\ind}{\mathbf 1}
\newcommand{\Var}{\operatorname{Var}}
\newcommand{\Cov}{\operatorname{Cov}}
\newcommand{\supp}{\operatorname{supp}}
\DeclareMathOperator*{\argmin}{arg\,min}
\DeclareMathOperator*{\argmax}{arg\,max}
\newcommand{\entrymeta}[3]{{\sffamily\small\color{muted}\textbf{\texttt{#1}}\quad|\quad #2\par #3\par}}
\newcommand{\Problemref}[1]{\texttt{#1}}
\newcommand{\JELref}[2]{\texttt{#2} (JEL \texttt{#1})}
\begin{document}
\section*{Recession and tail-risk forecasting}
\textbf{Problem ID:} \texttt{E37-594}\quad \textbf{JEL:} \texttt{E37}\par
% BEGIN ECONOMICS STATEMENT
\label{problem:OP-0239}
{\sffamily\small\color{muted}Original subject group: Business cycles and macro dynamics\par}
\label{definitions:problem:30007520}
\textbf{Audit classification:} Published research agenda. Real-time forecasting is an institutional agenda; a specific regret bound is editorial, not a published conjecture. Verification concerns the cited version, not first priority or proof that this exact editorial specification remains unsolved on October 8, 2026.
\subsection*{Definitions and precise research target}
\paragraph{Editorial mathematical specification.} Let \(X_t\) be a vector of real-time information available at date \(t\), including data vintages rather than later revisions. Define a recession-tail event \(D_{t,H}=1\{Y_{t+H}/Y_t-1<-q\}\), where \(Y_t\) is real output, \(H\) a fixed horizon, and \(q>0\) a predetermined threshold. A forecaster produces \(p_t=f_t(X_{\le t})\in[0,1]\). Specify a family \(\mathcal P\) of data-generating processes allowing at most \(K\) regime changes, bounded parameter drift between changes, and a minimum event frequency; these restrictions must be stated numerically before evaluation.
Construct a forecast rule whose excess expected proper-score loss is small relative to a complexity-restricted oracle comparator specified before evaluation. For Brier loss the target is a bound on
\[
\sup_{P\in\mathcal P}E_P\sum_{t=1}^T
\big[(p_t-D_{t,H})^2-(p_t^*-D_{t,H})^2\big],
\]
where \(p_t^*\) is the forecast supplied by a comparator in the specified complexity-restricted oracle class, evaluated with the same information available to the forecaster. The unrestricted conditional-probability oracle is only a lower-bound benchmark, not one against which small regret is guaranteed.
Evaluate simultaneous calibration and useful discrimination for rare downturns using strictly sequential holdouts spanning structural changes. Account for overlapping horizons and revised recession labels. Establish which information and stability restrictions make improvement possible; no uniformly accurate prediction is demanded for arbitrary unforeseeable regime changes. Report uncertainty rather than selecting a model on the same crises used to assess performance.

Require \(Y_t>0\) and choose \(0<q<1\) for a feasible negative-growth tail threshold. The outcome \(D_{t,H}\) becomes observable only at \(t+H\); the algorithm may not update from this label earlier. To make a regret statement possible, specify a conditional-probability model class of finite complexity, not only a bound on regime changes: for example a known finite set \(\{f_1,\ldots,f_N\}\) of real-time probability forecasts and an oracle index sequence with at most \(K\) switches. The benchmark loss is then that of the best such sequence, rather than an unrestricted oracle with unobservable information. Specify \(K=o(T)\), delayed feedback, and whether the guarantee concerns every outcome sequence or a restricted stochastic law. Calibration and discrimination are separate statistical properties and can conflict with uniform minimax requirements.
\subsection*{Variants and assumption changes}
The following are \emph{editorial variants}, not additional published open conjectures.
\begin{enumerate}
\item \emph{Finite switching benchmark.} Aggregate the \(N\) supplied forecasts under Brier loss and delayed labels; determine the attainable regret rate in \((T,N,K,H)\). This is a precisely defined methodological benchmark, not a claim of a newly unsolved online-learning theorem.
\item \emph{Conditional-density model.} Replace the finite pool by a declared smoothness class for the conditional GDP-growth density with quantified covariate dimension and variation budget. Study tail calibration and discrimination under that complexity restriction and the same real-time data vintages.
\end{enumerate}
\subsection*{References and source audit}
\begin{itemize}
\item European Central Bank. \emph{Forecasting and business cycle analysis}. Undated. Locator: Forecasting and related tools, selected areas of focus. Primary text checked. \url{https://www.ecb.europa.eu/pub/economic-research/research_agenda/forecasting/html/index.en.html}.
\end{itemize}
Real-time forecasting is an institutional agenda; a specific regret bound is editorial, not a published conjecture.
\begin{enumerate}\raggedright
\item\hypertarget{definitions:source:30007520:1}{} European Central Bank. Undated / date unverified. \emph{Forecasting and business cycle analysis}. Forecasting and related tools, selected areas of focus.
\url{https://www.ecb.europa.eu/pub/economic-research/research_agenda/forecasting/html/index.en.html}.
\end{enumerate}

\subsection*{Common conventions: Source mathematical conventions}

These conventions apply to each entry unless explicitly replaced.

\textbf{Models and identification.} A primitive model \(m\) specifies all probability laws, feasible choices, information, equilibrium and measurement rules used by its target. The admissible class \(\mathcal M\) is fixed independently of outcomes. If \(P_O(m)\) is its observable probability law and \(\tau(m)\) the target, its identified set at observable law \(P\) is
\[
 \mathcal I_\tau(P)=\{\tau(m):m\in\mathcal M,\ P_O(m)=P\}.
\]
Point identification means that this set is a singleton. An empty set signals incompatibility of data and assumptions. Coordinate bounds need not describe the joint set. Bounds are sharp only if every retained target is attainable or arbitrarily approachable by compatible models. Finite bounds require explicit restrictions on unobserved quantities; unrestricted confounding, hidden wealth, extrapolation or arbitrary equilibrium selection can yield uninformative sets.

\textbf{Causal interventions.} A potential outcome \(Y_i(p)\) is the outcome under a complete policy regime \(p\), including timing, financing, information and market spillovers. The target population and units are fixed before treatment. With interference, use the complete assignment vector or a justified exposure mapping; individual no-interference notation alone cannot identify an economy-wide intervention. A difference of mean potential outcomes does not require the cross-world joint distribution. Conditioning on post-treatment migration, survival, enrollment or employment changes the estimand and needs separate justification.

\textbf{Statistical statements.} An estimator, confidence set, forecast or test specifies a sampling law, model class and loss/coverage criterion. Uniform \((1-\alpha)\)-coverage means \(\inf_{m\in\mathcal M}\Pr_m\{\tau(m)\in C_n\}\geq1-\alpha\), or an explicitly asymptotic version. The whole parameter space trivially has coverage, so informativeness requires a stated width, risk or power objective. A forecasting improvement is relative to a named benchmark, information set and evaluation population.

\textbf{Equilibrium and optimization.} An equilibrium comprises feasible plans, optimality, beliefs where needed, budget constraints and all market/resource clearing. The primitives or a stated benchmark define each correspondence \(\mathcal E(p;m)\). Nonemptiness is assumed or proved, never inferred just from compact policy sets. A selected-equilibrium target uses a declared selection rule; a robust target quantifies over all permitted equilibria. Use suprema or infima unless attainment is established. An \(\varepsilon\)-optimal policy is feasible and within \(\varepsilon\) of the finite optimum value. Report an empty feasible set separately.

\textbf{Welfare and accounting.} Normative utility, interpersonal normalization, social weights, numeraire, discounting and the use of public receipts are primitives. Behavioral utility need not coincide with the welfare criterion. Household welfare includes ownership income and rents once; adding profits again to compensating variation generally double-counts them. A monetary equivalent is defined by an explicit transfer and price convention, with monotonicity and range conditions. Average effects, mechanism contributions, transfers and welfare losses are different targets. Infinite sums require convergence and dynamic feasible plans require terminal or no-Ponzi conditions.

\textbf{Source versions.} A working-paper date, revision date and journal publication year are distinguished. A DOI for a journal version is not automatically the DOI of an earlier manuscript. Locators refer to the stated version and printed pages where specified. A metadata-only check does not verify a claim about a full-text conclusion. Every entry's source subsection narrows the provenance claim accordingly.
% END ECONOMICS STATEMENT
\end{document}
